\chapter{Career Paths and Moves}\label{ch:in:career-paths-and-moves}

A trader who resigns at the start of February, a few days before the previous year's award is paid, may serve three
months of notice, spend three more on garden leave and then sit out a six-month non-compete, and forfeit on the way the
award for last year, the award accruing for this one and the deferred pay not yet vested. In this chapter's illustrative
package that is \$659\,167 forfeited and a year out of the market. A new employer's buyout is therefore part of the
offer, and the timing of a resignation is worth as much as a large raise. This chapter describes how people enter the
industry and move within it, puts a price on a move, and follows a career as a chain of states.

\section{Entry routes: internships, graduate programmes, lateral hires, doctorates}

\begin{definition}[Internship, return offer, graduate programme, lateral hire]\label{def:in:career-paths-and-moves:entry}
An \emph{internship}\index{internship} is a paid work placement of a few weeks or months for a student, designed by the
firm as a trial for a permanent job. A \emph{return offer}\index{return offer} is the offer of a permanent job made to an
intern at the end of the internship. A \emph{graduate programme}\index{graduate programme} is a firm's structured entry
route for new graduates, with training and often rotations before a permanent placement. A \emph{lateral
hire}\index{lateral hire} is an experienced person hired from another employer into a role at their level.
\end{definition}

For graduates the \omterm{def:in:career-paths-and-moves:entry}{internship} is the main door. The firms describe it as such: one market maker states that ``Our
\omterm{def:in:career-paths-and-moves:entry}{internship} program is built to prepare you for a full-time role''; another offers \omterm{def:in:career-paths-and-moves:entry}{internships} in quantitative trading,
research, machine learning and software engineering across its New York, London and Hong Kong offices. The hiring
process that leads to an \omterm{def:in:career-paths-and-moves:entry}{internship} or a graduate job (screening, assessments, interviews and the decision) is the
subject of Book~18.

\begin{dated}{2026-09}{How two firms describe their entry routes}\label{dat:in:career-paths-and-moves:entry}
Optiver: ``Our \omterm{def:in:career-paths-and-moves:entry}{internship} program is built to prepare you for a full-time role''; ``Our graduate roles are built around
learning through real work.'' Jane Street: \omterm{def:in:career-paths-and-moves:entry}{internships} in ``quantitative trading, machine learning, software engineering,
quantitative research, strategy and product'', across its New York, London and Hong Kong offices.
\end{dated}

A doctorate is an entry route of its own for research roles (chapter~17): the researcher joins after the doctorate,
sometimes after a research \omterm{def:in:career-paths-and-moves:entry}{internship} during it, at a level above a new graduate. \omterm{def:in:career-paths-and-moves:entry}{Lateral hires} fill the rest: an
experienced trader, researcher or engineer moving between firms, whose move this chapter prices. Chapter~29 follows the
pipelines from education.

\section{Bank to buy side and back}

Moves can run along any edge of the industry's map. From banks to the buy side: a \omterm{def:in:bank-quant:strat}{desk strategist} or \omterm{def:in:structurer-sales-and-sales-trader:structurer}{structurer}
(chapters~18 and~23) to a hedge fund, a trader to a platform (chapter~22). From the buy side to banks, for example when a
platform closes a pod or a fund shuts. Between buy-side firms, a researcher or a team moving with a strategy, which is
where non-competes and trade-secret law matter most (Book~16, chapter~11). Between technology firms and trading
(chapter~20), in both directions. No public data count these moves; what can be measured is their cost.

\begin{definition}[Notice period]\label{def:in:career-paths-and-moves:notice}
A \emph{notice period}\index{notice period} is the time between a resignation or dismissal and the end of employment,
set by the contract within the statutory minimum; the employee is paid and bound by the contract throughout, and may be
placed on garden leave for all or part of it.
\end{definition}

\begin{method}[The cost of a move]\label{met:in:career-paths-and-moves:move}
Let a job pay a base and an expected award $B$ paid in month $p$ for the previous year, a share $d$ deferred into $V$
annual tranches. In the steady state the unvested deferrals are worth $dB(V+1)/2$. Resigning in month $m$ forfeits them,
the previous year's award if $m\le p$, and the current year's accrual $B(m-1)/12$. Out of the market for $T$ months
(notice, garden leave, non-compete), the mover loses $BT/12$ of award, and the base for any unpaid months. A buyout
replaces a share $s$ of the forfeited awards with probability $q$. The move is neutral when a raise $\rho$ on base plus
award over $H$ years repays the net cost $C$: $\rho=C/((\text{base}+B)H)$.
\end{method}

\begin{example}[Move now or wait for the payment?]\label{ex:in:career-paths-and-moves:move}
Illustrative terms: base \$250\,000, expected award \$350\,000, 40\% deferred over three years, awards paid in February;
three months' notice, three of garden leave and a six-month paid non-compete; a buyout of all forfeited awards with
probability one-half; a three-year horizon. Unvested deferrals are \$280\,000. Resigning at the start of February
forfeits \$659\,167 and costs, net of the expected buyout, \$694\,167: a raise of 38.6\% on the \$600\,000 package. Waiting
one month, until the award is paid, cuts the \omterm{def:in:how-pay-works:rsu}{forfeiture} to \$338\,333 and the neutral raise to 30.5\%. Without a buyout
the February move needs 56.1\%; with a certain one, 21.1\% (\cref{fig:in:career-paths-and-moves:surface}).
\end{example}

\begin{omfigure}
\begin{tikzpicture}
\begin{axis}[width=10.5cm,height=5.6cm,xmin=0.5,xmax=12.5,ymin=0,ymax=60,xtick={1,2,3,4,5,6,7,8,9,10,11,12},
  xticklabels={Jan,Feb,Mar,Apr,May,Jun,Jul,Aug,Sep,Oct,Nov,Dec},xlabel={\small month of resignation},
  ylabel={\small neutral raise, \%},tick label style={font=\footnotesize},grid=major,grid style={gray!20},
  table/col sep=comma,legend style={legend cell align=left,font=\scriptsize,at={(0.5,-0.3)},anchor=north,legend columns=5,
  /tikz/every even column/.append style={column sep=6pt}}]
\addplot[black!60,thick,mark=o,mark options={solid}] table[x=month,y=q0]{figdata/industry/28-career-paths-and-moves/surface.csv};
\addplot[omThm,thick,mark=square*] table[x=month,y=q25]{figdata/industry/28-career-paths-and-moves/surface.csv};
\addplot[omDef,thick,mark=*] table[x=month,y=q50]{figdata/industry/28-career-paths-and-moves/surface.csv};
\addplot[omDef!50,thick,mark=triangle*] table[x=month,y=q75]{figdata/industry/28-career-paths-and-moves/surface.csv};
\addplot[black,thick,dashed,mark=diamond*,mark options={solid}] table[x=month,y=q100]{figdata/industry/28-career-paths-and-moves/surface.csv};
\legend{no buyout,$q=0.25$,$q=0.5$,$q=0.75$,certain buyout}
\end{axis}
\end{tikzpicture}

\omcaption{The raise on base plus award that makes the chapter's illustrative move neutral over three years, by month of
resignation and probability $q$ of a full buyout of forfeited awards. The award is paid in February: resigning before
the payment forfeits it, and the year's accrual grows each month after. Data: \texttt{in\_career.surface}, from
\texttt{firm.careerpath.breakeven\_raise}.}\label{fig:in:career-paths-and-moves:surface}
\end{omfigure}

The shape of \cref{fig:in:career-paths-and-moves:surface} is the practical lesson. The cost falls sharply once an award is
paid and then rises through the year as the next one accrues; a buyout flattens the payment-day cliff, because it
replaces forfeited awards but not the accrual or the time out. The year out of the market costs \$350\,000 of award in
every month: the non-compete is the largest single item.

\section{Non-competes and garden leave in practice}

Book~16, chapter~11, sets out the law: non-compete clauses, garden leave, non-solicitation and confidentiality, and how
courts treat them. The practice depends on where the job is.

\begin{dated}{2026-09}{Notice and non-compete law}\label{dat:in:career-paths-and-moves:law}
United Kingdom: statutory minimum notice from the employer of one week a year of service up to twelve weeks, and from
the employee one week (Employment Rights Act 1996, s.~86); a contract may set longer periods. The government's
working paper of 26 November 2025 consulted, until 18 February 2026, on options including a statutory limit on the
length of non-competes, a ban, or a ban below a salary threshold; the previous government had announced a three-month
limit on 10 May 2023 that was not enacted. United States: the FTC's non-compete rule was stopped by a district court on
20 August 2024, and on 5 September 2025 the Commission voted to dismiss its appeals and accede to the vacatur; the
rule is not in force. California: a non-compete void under its law ``is unenforceable regardless of where and when the
contract was signed'' (Business and Professions Code s.~16600.5).
\end{dated}

The employer's side of the calculation is Book~16's \texttt{firm.gardenleave}: paying someone to stay out of the market is
worth it while the strategy knowledge they carry decays more slowly than the pay costs. With the chapter's illustrative
inputs (a strategy earning \$5 million a year, 20\% of it lost to a competitor who learns it, a one-year half-life of the
edge, and \$600\,000 a year of pay during the leave) the best length is 0.74 years, about nine months, which is the order of
the leaves the example assumes. A firm that pays for a leave is buying time for its strategies to change.

\section{Going independent and starting a firm}

Some moves are out of employment altogether: a \omterm{def:in:portfolio-manager-and-pod-analyst:roles}{portfolio manager} launches a fund, a group of traders starts a
proprietary firm, an engineer starts a vendor. Book~16, chapter~25, covers launching a fund and raising capital, with the
emerging-manager programmes and seed investors that back new funds, and the portability of a track record that decides
whether investors believe it. Chapter~3's lineage of trading firms records how several options market makers began on
exchange floors; it has no sourced \omterm{def:in:the-options-market-making-houses:spinout}{spin-out} edge, because the public record rarely documents who left which firm to
found another. The move's economics are those of the partner in chapter~25 without the salary: capital at risk from the
first day and income only from profits.

\section{Exits}

Careers also leave the industry: to technology firms, to academia, to the official sector (regulators and central banks,
chapter~9), to corporate treasury and risk, or to retirement after a short career. The illustrative career chain of
\cref{fig:in:career-paths-and-moves:fan} gives each year a 5\% chance of leaving; with its other assumptions, 36.7\% of
graduates have left by year ten and 21.5\% run a portfolio. Its pay paths spread widely: the median rises to a senior
researcher's pay by year eight and falls back as exits accumulate, while the 90th percentile reaches a \omterm{def:in:portfolio-manager-and-pod-analyst:roles}{portfolio
manager}'s from year seven.

\begin{omfigure}
\begin{tikzpicture}
\begin{axis}[width=10.5cm,height=5.6cm,xmin=1,xmax=15,ymin=0,ymax=1800,xlabel={\small year of career},
  ylabel={\small annual pay, \$ thousand},tick label style={font=\footnotesize},grid=major,grid style={gray!20},
  table/col sep=comma,legend style={font=\scriptsize,at={(0.03,0.97)},anchor=north west,legend cell align=left}]
\addplot[name path=hi,draw=none,forget plot] table[x=year,y=p90]{figdata/industry/28-career-paths-and-moves/fan.csv};
\addplot[name path=lo,draw=none,forget plot] table[x=year,y=p10]{figdata/industry/28-career-paths-and-moves/fan.csv};
\addplot[omDef!20] fill between[of=hi and lo];
\addplot[omDef,thick] table[x=year,y=mean]{figdata/industry/28-career-paths-and-moves/fan.csv};
\addplot[omThm,thick,dashed] table[x=year,y=p50]{figdata/industry/28-career-paths-and-moves/fan.csv};
\legend{10th to 90th percentile,mean,median}
\end{axis}
\end{tikzpicture}

\omcaption{Annual pay along an illustrative career chain (graduate, researcher, senior researcher, \omterm{def:in:portfolio-manager-and-pod-analyst:roles}{portfolio manager},
left the industry) with labelled transition probabilities and pay by state: mean, median and 10th--90th percentile band
over 20\,000 careers. Every number is an assumption; the shape, a widening fan, is the point. Data:
\texttt{in\_career.career}, from \texttt{firm.careerpath.simulate}.}\label{fig:in:career-paths-and-moves:fan}
\end{omfigure}

\section{Tutorial: move now or wait for the vest?}\label{tut:in:career-paths-and-moves:tut}

\begin{tutorial}
\textbf{Goal.} Price a move by its month and its buyout, and follow a career as a Markov chain.
\textbf{End state:} \cref{fig:in:career-paths-and-moves:surface,fig:in:career-paths-and-moves:fan}.
\begin{tutsteps}
\item \textbf{The package.} \texttt{firm.careerpath.Job(base, award, deferral, vest\_years, pay\_month)}; \texttt{forfeited}
gives the three forfeited items by month.
\item \textbf{The move.} \texttt{move\_cost} and \texttt{breakeven\_raise} (\cref{lst:in:career-paths-and-moves:move}).
\omcode{code/firm/careerpath/firm_careerpath.py}{50}{62}{The net cost of a move and the raise that makes it neutral.}\label{lst:in:career-paths-and-moves:move}
\item \textbf{The employer's view.} Book~16's \texttt{firm.gardenleave.best\_length}.
\item \textbf{The career.} \texttt{Chain(states, P, pay)}, \texttt{simulate} and \texttt{pay\_paths}; the lineage graph's
\omterm{def:in:the-options-market-making-houses:spinout}{spin-outs} through \texttt{spinouts}.
\end{tutsteps}
\end{tutorial}

The surface's corners: resigning at the start of February needs 56.1\% without a buyout and 21.1\% with a certain one;
resigning at the start of December, 52.8\% and 37.3\%.

\textbf{What to change next.} Discount the raise and let the award vary; let the buyout cover only deferred pay; add a
good-leaver clause (chapter~13); estimate a chain's transition probabilities from public profiles, with their bias.

\section{Build: firm.careerpath}\label{bld:in:career-paths-and-moves:careerpath}

\begin{build}
\bfield{Purpose} Price a move between employers, and model a career as a chain of states with labelled assumptions.
\bfield{Interface} \texttt{firm.careerpath}: \texttt{Job}; \texttt{unvested}; \texttt{forfeited(job, month)};
\texttt{move\_cost(job, month, notice\_m, garden\_m, noncompete\_m, paid\_noncompete, q, s)};
\texttt{breakeven\_raise(job, month, horizon\_y, \dots)}; \texttt{Chain(states, P, pay)}; \texttt{simulate(chain, start,
years, n, rng)}; \texttt{pay\_paths}; \texttt{spinouts(graph)}; with \texttt{firm.gardenleave} and \texttt{firm.lineage}.
\bfield{Rules} Bad-leaver \omterm{def:in:how-pay-works:rsu}{forfeiture}; no award accrues out of the market; rows of the chain sum to one or the simulation
refuses; every parameter labelled by the caller.
\bfield{Acceptance tests} \texttt{code/firm/careerpath/tests/}: the unpaid award is forfeited only before its payment
month; the worst month is the one before payment; a certain buyout offsets the forfeited awards; a deterministic chain
moves as specified.
\bfield{Stretch} Discounting and award uncertainty (with \texttt{firm.payoffer}); a chain with duration-dependent
promotion; team moves.
\end{build}

\begin{omsources}
\item Employment Rights Act 1996, s.~86; UK Department for Business and Trade (2025), working paper on non-compete reform.
\item Federal Trade Commission (2025), press release on the non-compete rule; California Business and Professions Code
s.~16600.5.
\item Optiver and Jane Street, students and internships pages.
\item Book~16, chapters~11 and~25; Book~18 on hiring.
\end{omsources}

\section{Exercises}

\begin{exercise}[$\star$]\label{exo:in:career-paths-and-moves:1}
With the example's package, what are the unvested deferrals, and what does resigning at the start of March forfeit?
\end{exercise}

\begin{exercise}[$\star$]\label{exo:in:career-paths-and-moves:2}
What is the UK statutory minimum notice for an employer to give someone with seven years of service?
\end{exercise}

\begin{exercise}[$\star$]\label{exo:in:career-paths-and-moves:3}
How much award does a year out of the market cost in the example?
\end{exercise}

\begin{exercise}[$\star\star$]\label{exo:in:career-paths-and-moves:4}
A certain buyout makes a February move cheaper than a March move without one. Why is it still not free?
\end{exercise}

\begin{exercise}[$\star\star$]\label{exo:in:career-paths-and-moves:5}
A new employer offers a 25\% raise and a 50\% chance of a buyout. In which months is the move worth it?
\end{exercise}

\begin{exercise}[$\star\star$]\label{exo:in:career-paths-and-moves:6}
Why would a candidate care whether a non-compete is paid?
\end{exercise}

\begin{exercise}[$\star\star\star$]\label{exo:in:career-paths-and-moves:7}
\emph{Coding.} Make the non-compete unpaid. What is the neutral raise for a February resignation with a 50\% buyout
chance?
\end{exercise}

\begin{exercise}[$\star\star\star$]\label{exo:in:career-paths-and-moves:8}
\emph{Find the flaw.} ``The model shows that a fifth of graduates become \omterm{def:in:portfolio-manager-and-pod-analyst:roles}{portfolio managers} within ten years.''
\end{exercise}

\section{Problem: Move Now or Wait for the Vest?}

\begin{problem}[{Weekend problem --- move now or wait for the vest?}]\label{pb:in:career-paths-and-moves:1}
A trader on the example's package receives an offer in late January and asks when to resign.

\medskip\noindent\textbf{Part I --- The routes.}
\begin{enumerate}
\item Define an \omterm{def:in:career-paths-and-moves:entry}{internship}, a \omterm{def:in:career-paths-and-moves:entry}{return offer}, a \omterm{def:in:career-paths-and-moves:entry}{graduate programme} and a \omterm{def:in:career-paths-and-moves:entry}{lateral hire}.
\item How do the firms describe their \omterm{def:in:career-paths-and-moves:entry}{internships} and graduate roles?
\item Where does a doctorate lead?
\item Name the common edges of moves between employers.
\item What public data count them?
\end{enumerate}

\medskip\noindent\textbf{Part II --- The cost.}
\begin{enumerate}[resume]
\item Define the \omterm{def:in:career-paths-and-moves:notice}{notice period} and state the UK statutory minimum.
\item State the method.
\item Compute the \omterm{def:in:how-pay-works:rsu}{forfeiture} and the net cost for a February and a March resignation.
\item Compute the neutral raise for each with a 50\% buyout chance.
\item What is the largest single item, and why?
\end{enumerate}

\medskip\noindent\textbf{Part III --- The law and the employer.}
\begin{enumerate}[resume]
\item What is the status of the FTC's non-compete rule?
\item What has the UK government proposed, and when?
\item What does California's s.~16600.5 add?
\item What garden-leave length does the employer's model give, and from what?
\item What does a \omterm{def:in:the-options-market-making-houses:spinout}{spin-out} need that a move does not?
\end{enumerate}

\medskip\noindent\textbf{Part IV --- The verdict.}
\begin{enumerate}[resume]
\item State the \emph{named result}: the raise that makes the move neutral as a function of the month of leaving and the
buyout probability, and its value on the day before the award is paid.
\item What should he ask the new employer to buy out?
\item How would a good-leaver clause change the answer?
\item What does the career chain say, and not say?
\item In two sentences, advise him.
\end{enumerate}
\end{problem}

\section{Interview questions}

\begin{interviewq}[$\star$ \iqroles{trader, researcher}]\label{iq:in:career-paths-and-moves:1}
Why do you want to leave your current firm?
\end{interviewq}

\begin{interviewq}[$\star$ \iqroles{researcher}]\label{iq:in:career-paths-and-moves:2}
What can you tell us about your current strategies, and what can you not?
\end{interviewq}

\begin{interviewq}[$\star\star$ \iqroles{trader}]\label{iq:in:career-paths-and-moves:3}
You have a six-month non-compete. How would you use the time?
\end{interviewq}

\begin{interviewq}[$\star\star$ \iqroles{researcher, trader}]\label{iq:in:career-paths-and-moves:4}
How would you show us your track record without disclosing your employer's information?
\end{interviewq}

\begin{interviewq}[$\star\star$ \iqroles{bank, researcher}]\label{iq:in:career-paths-and-moves:5}
Why move from a bank to a fund, and what will you miss?
\end{interviewq}

\begin{interviewq}[$\star\star\star$ \iqroles{trader, researcher}]\label{iq:in:career-paths-and-moves:6}
Price the deferred pay you would forfeit by joining us, and propose a buyout.
\end{interviewq}
